\chapter{Delayed Ejaculation}
\index{Delayed Ejaculation}

\section*{Definition and Overview}
Delayed ejaculation is a sexual difficulty in which a man has persistent or recurrent difficulty achieving ejaculation, or an unusually long delay before ejaculation occurs, despite adequate stimulation and the desire to ejaculate. It can occur in all sexual situations or only in specific ones, and it ranges from mild delay to the complete inability to ejaculate. Delayed ejaculation can cause distress and relationship difficulties. It is less common than premature ejaculation but is a recognised and treatable condition with psychological and medical causes.

\section*{Epidemiology}
Delayed ejaculation is the least common of the male sexual dysfunctions, affecting around 1--5\% of men, with prevalence increasing with age. It is more common in men over 50 and in those with medical conditions such as diabetes, neurological disease, or prostate surgery. Medications, particularly antidepressants and some blood pressure drugs, are an important cause. Many men experience situational delay at some point, but the condition is considered a disorder when it is persistent, distressing, or causes relationship problems.

\section*{Aetiology and Pathophysiology}
\begin{itemize}
    \item \textbf{Psychological factors:} anxiety, stress, depression, performance concerns, relationship conflict, and religious or cultural beliefs about sex
    \item \textbf{Medications:} antidepressants (especially SSRIs and SNRIs), some antipsychotics, blood pressure medicines, and opioid painkillers
    \item \textbf{Medical conditions:} diabetes with nerve damage, multiple sclerosis, spinal cord injury, and hypothyroidism
    \item \textbf{Pelvic surgery:} prostatectomy, bladder surgery, and other pelvic procedures can affect ejaculatory function
    \item \textbf{Substance use:} alcohol and recreational drugs can impair ejaculation
    \item \textbf{Neurological and hormonal factors:} reduced sensitivity, ageing changes, and low testosterone
    \item \textbf{Reduced penile sensation:} from circumcision-related scarring, neuropathy, or excessive masturbation technique
\end{itemize}

\section*{Clinical Features}
\begin{itemize}
    \item Difficulty ejaculating despite prolonged stimulation and sexual arousal
    \item Delay that is out of proportion to what the man or partner considers normal
    \item Situational patterns: for example, ability to ejaculate with masturbation but not with a partner, or vice versa
    \item Reduced sexual satisfaction, frustration, and avoidance of sexual activity
    \item Relationship tension and reduced intimacy
    \item Absence of ejaculation during orgasm (anejaculation) in the most severe form
    \item Infertility when ejaculation is consistently absent in partnered sex
\end{itemize}

\section*{Differential Diagnosis}
\begin{itemize}
    \item Premature ejaculation, the opposite pattern, which is far more common
    \item Erectile dysfunction, which may coexist and complicate the picture
    \item Low sexual desire, which reduces arousal and makes ejaculation difficult
    \item Retrograde ejaculation, where semen passes into the bladder instead of out through the penis
    \item Obstructive causes of ejaculatory failure, such as prostate or duct problems
    \item Medication effects, which should always be considered
\end{itemize}

\section*{When to Seek Medical Care}
Men who are distressed by delayed ejaculation, or whose partners are affected, should discuss it with a doctor, as it is a legitimate health concern with effective treatments. A GP can review medications, check for underlying conditions, and refer to a sexual health physician, urologist, or psychologist. Men with diabetes, prostate surgery, or neurological conditions who develop ejaculatory difficulty should also seek review.

\textbf{Contact your local emergency services for:}
\begin{itemize}
    \item Sudden loss of ejaculatory function with back or pelvic pain, weakness, or numbness, which may indicate a neurological emergency
\end{itemize}

\section*{Diagnosis and Investigations}
\begin{itemize}
    \item \textbf{Detailed history:} onset, situations, relationship factors, medication use, and associated sexual problems
    \item \textbf{Medication review:} identifying drugs that affect ejaculation, such as antidepressants
    \item \textbf{Physical examination:} neurological, urological, and endocrine assessment
    \item \textbf{Blood tests:} testosterone, thyroid function, and blood glucose where indicated
    \item \textbf{Specialist tests:} urological investigation for retrograde ejaculation or obstruction, and assessment of penile sensation
    \item \textbf{Psychological assessment:} exploration of anxiety, mood, and relationship factors
\end{itemize}

\section*{Management}
\begin{itemize}
    \item \textbf{Addressing medication:} adjusting, switching, or stopping culprit medications under medical supervision
    \item \textbf{Psychological therapy:} cognitive-behavioural therapy and sex therapy address anxiety, performance concerns, and relationship issues
    \item \textbf{Sensate focus:} structured exercises to rebuild intimacy and reduce pressure around ejaculation
    \item \textbf{Stimulation techniques:} adjusting masturbation patterns and trying different stimulation, including vibrators, can help
    \item \textbf{Treating underlying conditions:} diabetes control, thyroid treatment, and management of neurological conditions
    \item \textbf{Medication options:} in some cases, drugs that may facilitate ejaculation are used under specialist guidance
    \item \textbf{Assisted reproduction:} for couples wanting children when ejaculation remains absent, sperm retrieval techniques are available
    \item \textbf{Couple counselling:} support for both partners to reduce pressure and improve communication
\end{itemize}

\section*{Complications}
\begin{itemize}
    \item Distress, anxiety, and reduced self-esteem
    \item Relationship conflict, reduced intimacy, and avoidance of sex
    \item Difficulty conceiving when ejaculation is absent
    \item Depression and reduced quality of life
    \item Effects on the partner's satisfaction and wellbeing
\end{itemize}

\section*{Prognosis and Outcomes}
The outlook for delayed ejaculation is generally good once the cause is identified. When it is medication-related, adjusting treatment often restores function. Psychological causes respond well to sex therapy and cognitive-behavioural approaches, particularly when the problem is situational. Underlying medical conditions can often be managed, although ejaculatory function may not fully recover after extensive prostate surgery. Most men and couples find that open communication and professional support significantly improve their sexual relationship and satisfaction.

\section*{Prevention and Self-Care}
\begin{itemize}
    \item Talk openly with your partner about sex, expectations, and concerns
    \item Discuss possible sexual side effects before starting new medications
    \item Reduce performance pressure by focusing on intimacy and pleasure rather than ejaculation as a goal
    \item Limit alcohol and avoid recreational drugs before sex
    \item Manage stress, sleep, and overall health
    \item Seek help early rather than avoiding sexual activity
    \item Consider professional support from a GP, sexual health service, or psychologist
\end{itemize}

\section*{Sources and Further Reading}
The following reputable sources provide further information for healthcare students and patients seeking reliable, current advice about delayed ejaculation.

\begin{itemize}
    \item NHS. \textit{Delayed ejaculation: causes and treatment.} https://www.nhs.uk/conditions/sexual-problems-in-men/
    \item Mayo Clinic. \textit{Delayed ejaculation: symptoms and treatment.} https://www.mayoclinic.org/diseases-conditions/delayed-ejaculation/
    \item MSD Manual. \textit{Delayed ejaculation: professional overview.} https://www.msdmanuals.com/
    \item MedlinePlus. \textit{Ejaculation problems.} https://medlineplus.gov/
    \item Healthdirect Australia. \textit{Sexual problems in men: information and support.} https://www.healthdirect.gov.au/
    \item Andrology Australia. \textit{Male sexual health information.} https://www.andrologyaustralia.org/
\end{itemize}
